% =============================================================================
% chapters/ch03_theory.tex — Theory & Hypotheses
% Pulled into manuscript.tex via % =============================================================================
% chapters/ch03_theory.tex — Theory & Hypotheses
% Pulled into manuscript.tex via % =============================================================================
% chapters/ch03_theory.tex — Theory & Hypotheses
% Pulled into manuscript.tex via % =============================================================================
% chapters/ch03_theory.tex — Theory & Hypotheses
% Pulled into manuscript.tex via \include{chapters/ch03_theory}.
% =============================================================================
\chapter{Theory: Aid Races and Aid Withdrawal}
\label{ch:theory}

\textit{Placeholder.} Develop the core theory: when a donor faces a strategic
rival entering a recipient it values, it either escalates (an \emph{aid race})
or disengages (\emph{aid withdrawal}). Which response dominates depends on
whether the recipient lies within a contested sphere of influence and on the
type of good the aid delivers.

\section{Spheres of influence, operationalized}
\label{sec:soi-operational}
The first-pass operationalization used continent as a stand-in for ``sphere of
influence.'' That was always a coarse proxy -- continent doesn't distinguish a
country firmly in the US camp from a genuinely contested neighbor -- and
Chapter~\ref{ch:quant} shows the region cut produces a diffuse, largely
non-significant pattern once standard errors are handled correctly. The
refined operationalization (voting-alignment component only; geographic
proximity and historical-ties components are not yet built) uses each
country's UN General Assembly ideal point (Voeten, Sep.\ 2023 release) to
locate it relative to both powers:
\[
\text{rel\_align}_{i,t} = |\text{ideal}_{i,t} - \text{ideal}_{\text{CHN},t}| - |\text{ideal}_{i,t} - \text{ideal}_{\text{US},t}|
\]
Large positive values mean a country votes far from China and close to the US
(\emph{US-aligned}); large negative values mean the reverse (\emph{China-aligned});
values near zero mean the country is roughly equidistant --
genuinely \emph{contested}. Countries are terciled into these three blocs
\textbf{within year} (ideal-point scales are not necessarily comparable across
UNGA sessions) \textbf{and within the US-aid-recipient sample} (not the full
UN membership, so a bloc means ``relative to other aid recipients''), using
the bloc as of $t-1$ so this year's aid decision cannot
retroactively define the very bloc it is tested against. This SOI measure is
tested directly in Chapter~\ref{ch:quant}.

\section{H1: Aid races vs.\ aid withdrawal}
\textbf{H1.} In contested spheres, US and Chinese aid co-move positively (a
race); where the US has effectively already lost a recipient to China, US aid
falls as Chinese aid rises (withdrawal) -- a country already secured in the
US's own camp carries no sharp prediction either way. The prospectus finding
of $+0.113$ in Africa and negative correlations elsewhere was the original
motivating pattern; Chapter~\ref{ch:quant} tests this directly against the
alignment-based SOI blocs above, and the result is not what H1 predicts.

\section{H2: Private vs.\ public goods delivery}
\textit{Placeholder.} \textbf{H2.} The race is sharper where aid is delivered
as private/excludable goods (e.g., government-to-government projects) than where
it flows through NGOs to broad public goods. Tested in
\texttt{analyses/02\_h2\_private\_goods/}.
.
% =============================================================================
\chapter{Theory: Aid Races and Aid Withdrawal}
\label{ch:theory}

\textit{Placeholder.} Develop the core theory: when a donor faces a strategic
rival entering a recipient it values, it either escalates (an \emph{aid race})
or disengages (\emph{aid withdrawal}). Which response dominates depends on
whether the recipient lies within a contested sphere of influence and on the
type of good the aid delivers.

\section{Spheres of influence, operationalized}
\label{sec:soi-operational}
The first-pass operationalization used continent as a stand-in for ``sphere of
influence.'' That was always a coarse proxy -- continent doesn't distinguish a
country firmly in the US camp from a genuinely contested neighbor -- and
Chapter~\ref{ch:quant} shows the region cut produces a diffuse, largely
non-significant pattern once standard errors are handled correctly. The
refined operationalization (voting-alignment component only; geographic
proximity and historical-ties components are not yet built) uses each
country's UN General Assembly ideal point (Voeten, Sep.\ 2023 release) to
locate it relative to both powers:
\[
\text{rel\_align}_{i,t} = |\text{ideal}_{i,t} - \text{ideal}_{\text{CHN},t}| - |\text{ideal}_{i,t} - \text{ideal}_{\text{US},t}|
\]
Large positive values mean a country votes far from China and close to the US
(\emph{US-aligned}); large negative values mean the reverse (\emph{China-aligned});
values near zero mean the country is roughly equidistant --
genuinely \emph{contested}. Countries are terciled into these three blocs
\textbf{within year} (ideal-point scales are not necessarily comparable across
UNGA sessions) \textbf{and within the US-aid-recipient sample} (not the full
UN membership, so a bloc means ``relative to other aid recipients''), using
the bloc as of $t-1$ so this year's aid decision cannot
retroactively define the very bloc it is tested against. This SOI measure is
tested directly in Chapter~\ref{ch:quant}.

\section{H1: Aid races vs.\ aid withdrawal}
\textbf{H1.} In contested spheres, US and Chinese aid co-move positively (a
race); where the US has effectively already lost a recipient to China, US aid
falls as Chinese aid rises (withdrawal) -- a country already secured in the
US's own camp carries no sharp prediction either way. The prospectus finding
of $+0.113$ in Africa and negative correlations elsewhere was the original
motivating pattern; Chapter~\ref{ch:quant} tests this directly against the
alignment-based SOI blocs above, and the result is not what H1 predicts.

\section{H2: Private vs.\ public goods delivery}
\textit{Placeholder.} \textbf{H2.} The race is sharper where aid is delivered
as private/excludable goods (e.g., government-to-government projects) than where
it flows through NGOs to broad public goods. Tested in
\texttt{analyses/02\_h2\_private\_goods/}.
.
% =============================================================================
\chapter{Theory: Aid Races and Aid Withdrawal}
\label{ch:theory}

\textit{Placeholder.} Develop the core theory: when a donor faces a strategic
rival entering a recipient it values, it either escalates (an \emph{aid race})
or disengages (\emph{aid withdrawal}). Which response dominates depends on
whether the recipient lies within a contested sphere of influence and on the
type of good the aid delivers.

\section{Spheres of influence, operationalized}
\label{sec:soi-operational}
The first-pass operationalization used continent as a stand-in for ``sphere of
influence.'' That was always a coarse proxy -- continent doesn't distinguish a
country firmly in the US camp from a genuinely contested neighbor -- and
Chapter~\ref{ch:quant} shows the region cut produces a diffuse, largely
non-significant pattern once standard errors are handled correctly. The
refined operationalization (voting-alignment component only; geographic
proximity and historical-ties components are not yet built) uses each
country's UN General Assembly ideal point (Voeten, Sep.\ 2023 release) to
locate it relative to both powers:
\[
\text{rel\_align}_{i,t} = |\text{ideal}_{i,t} - \text{ideal}_{\text{CHN},t}| - |\text{ideal}_{i,t} - \text{ideal}_{\text{US},t}|
\]
Large positive values mean a country votes far from China and close to the US
(\emph{US-aligned}); large negative values mean the reverse (\emph{China-aligned});
values near zero mean the country is roughly equidistant --
genuinely \emph{contested}. Countries are terciled into these three blocs
\textbf{within year} (ideal-point scales are not necessarily comparable across
UNGA sessions) \textbf{and within the US-aid-recipient sample} (not the full
UN membership, so a bloc means ``relative to other aid recipients''), using
the bloc as of $t-1$ so this year's aid decision cannot
retroactively define the very bloc it is tested against. This SOI measure is
tested directly in Chapter~\ref{ch:quant}.

\section{H1: Aid races vs.\ aid withdrawal}
\textbf{H1.} In contested spheres, US and Chinese aid co-move positively (a
race); where the US has effectively already lost a recipient to China, US aid
falls as Chinese aid rises (withdrawal) -- a country already secured in the
US's own camp carries no sharp prediction either way. The prospectus finding
of $+0.113$ in Africa and negative correlations elsewhere was the original
motivating pattern; Chapter~\ref{ch:quant} tests this directly against the
alignment-based SOI blocs above, and the result is not what H1 predicts.

\section{H2: Private vs.\ public goods delivery}
\textit{Placeholder.} \textbf{H2.} The race is sharper where aid is delivered
as private/excludable goods (e.g., government-to-government projects) than where
it flows through NGOs to broad public goods. Tested in
\texttt{analyses/02\_h2\_private\_goods/}.
.
% =============================================================================
\chapter{Theory: Aid Races and Aid Withdrawal}
\label{ch:theory}

\textit{Placeholder.} Develop the core theory: when a donor faces a strategic
rival entering a recipient it values, it either escalates (an \emph{aid race})
or disengages (\emph{aid withdrawal}). Which response dominates depends on
whether the recipient lies within a contested sphere of influence and on the
type of good the aid delivers.

\section{Spheres of influence, operationalized}
\label{sec:soi-operational}
The first-pass operationalization used continent as a stand-in for ``sphere of
influence.'' That was always a coarse proxy -- continent doesn't distinguish a
country firmly in the US camp from a genuinely contested neighbor -- and
Chapter~\ref{ch:quant} shows the region cut produces a diffuse, largely
non-significant pattern once standard errors are handled correctly. The
refined operationalization (voting-alignment component only; geographic
proximity and historical-ties components are not yet built) uses each
country's UN General Assembly ideal point (Voeten, Sep.\ 2023 release) to
locate it relative to both powers:
\[
\text{rel\_align}_{i,t} = |\text{ideal}_{i,t} - \text{ideal}_{\text{CHN},t}| - |\text{ideal}_{i,t} - \text{ideal}_{\text{US},t}|
\]
Large positive values mean a country votes far from China and close to the US
(\emph{US-aligned}); large negative values mean the reverse (\emph{China-aligned});
values near zero mean the country is roughly equidistant --
genuinely \emph{contested}. Countries are terciled into these three blocs
\textbf{within year} (ideal-point scales are not necessarily comparable across
UNGA sessions) \textbf{and within the US-aid-recipient sample} (not the full
UN membership, so a bloc means ``relative to other aid recipients''), using
the bloc as of $t-1$ so this year's aid decision cannot
retroactively define the very bloc it is tested against. This SOI measure is
tested directly in Chapter~\ref{ch:quant}.

\section{H1: Aid races vs.\ aid withdrawal}
\textbf{H1.} In contested spheres, US and Chinese aid co-move positively (a
race); where the US has effectively already lost a recipient to China, US aid
falls as Chinese aid rises (withdrawal) -- a country already secured in the
US's own camp carries no sharp prediction either way. The prospectus finding
of $+0.113$ in Africa and negative correlations elsewhere was the original
motivating pattern; Chapter~\ref{ch:quant} tests this directly against the
alignment-based SOI blocs above, and the result is not what H1 predicts.

\section{H2: Private vs.\ public goods delivery}
\textit{Placeholder.} \textbf{H2.} The race is sharper where aid is delivered
as private/excludable goods (e.g., government-to-government projects) than where
it flows through NGOs to broad public goods. Tested in
\texttt{analyses/02\_h2\_private\_goods/}.
